% SGA 6, Exposición VI: fragmento español alineado con la fuente.
% Alcance: proposición 1.3 y demostración, líneas 8588--8624 inclusive.
% Fuente francesa: sga6_fr_r1/src/fr.tex.
% SHA-256 de la fuente francesa completa:
% 0D2158B9FE3B466309A92165727A0B191422E05FBAA0DACD5806E8322C54BA1A.
\providecommand{\sisi}[2]{#2}

\textbf{Proposición 1.3.} Sean \(S\) un esquema cuasicompacto, \(E\) un \(\mathcal O_S\)-módulo localmente libre de tipo finito, y \(X\) el fibrado proyectivo asociado. Sea \(F\) un \(\mathcal O_X\)-módulo de presentación finita, plano sobre \(S\). Existe un entero \(n_0\) tal que, para \(n\geq n_0\), se tiene:
\[
\begin{array}{ll}
\text{i)}& R^i f_*(F(n))=0\quad\text{para }i\geq 1;\\[2mm]
\text{ii)}& f_*(F(n))\text{ es localmente libre de tipo finito sobre }S;\\[2mm]
\text{iii)}& M=\displaystyle\bigoplus_{n\geq n_0} f_*(F(n))\text{ es un }\Sym(E)\text{-módulo graduado de presentación finita (1.2).}
\end{array}
\]

Las propiedades buscadas son de naturaleza local sobre \(S\); como \(S\) es cuasicompacto, se puede, tomando el mayor de los enteros encontrados, reducir al caso en que \(S\) es afín.

Supongamos primero que \(S\) es noetheriano.

La propiedad i) resulta entonces del teorema de Serre.

La propiedad ii) resulta de EGA III 7.9.10, eligiendo \(n_0\) de manera que se verifique la propiedad i).

Según EGA II 3.4.5, existe un \(\Sym(E)\)-módulo graduado cuasicoherente \(M'\), de tipo finito, tal que \(F=\Projz(M')\). Por otra parte, \(\Gamma_*(F)=\Gamma_*(\Projz(M'))=\displaystyle\bigoplus_{n\in\mathbb Z} f_*(F(n))\) es TN-isomorfo a \(M'\) según \sisi{ECA}{EGA} III 2.3.1; existe, pues, un entero \(n_1\) tal que \(\displaystyle\bigoplus_{n\geq n_1} f_*(F(n))\simeq \bigoplus_{n\geq n_1} M'_n\),
de modo que \(\bigoplus_{n\geq n_1} f_*(F(n))\) es un \(\Sym(E)\)-módulo de tipo finito, y por tanto de presentación finita bajo las hipótesis noetherianas; por lo demás, el resultado es válido para todo entero \(n_1\), pues los \(f_*(F(n))\) son coherentes según el teorema de Serre.

La proposición queda, pues, demostrada cuando \(S\) es noetheriano.

Supongamos ahora que se tiene un cuadrado cartesiano:
\[
\begin{tikzcd}[column sep=large,row sep=large]
X \arrow[r,"g"] \arrow[d,"f"'] & X_0 \arrow[d,"f_0"]\\
S \arrow[r,"g_0"] & S_0,
\end{tikzcd}
\]
con \(S_0\) noetheriano y \(X_0=\mathbf P(E_0)\); sea \(F_0\) un \(\sisi{\mathcal O_X}{\mathcal O_{X_0}}\)-módulo de presentación finita, plano sobre \(S_0\), y sea \(F=g^*(F_0)\). Por lo que precede, la proposición es verdadera para \(F_0\). Según EGA III 7.7.5 y 7.8.5, el teorema de cambio de base es verdadero bajo las hipótesis anteriores; se deduce inmediatamente que el teorema es verdadero para \(F\).

Para concluir la demostración, basta, pues, mostrar que el caso general (con \(S\) afín) se puede obtener a partir del caso noetheriano mediante cambio de base.

Como \(S\) es afín, \(S\) es límite proyectivo filtrante de esquemas afines noetherianos (si \(S=\Spec(A)\), se escribe \(A\) como límite inductivo de sus subálgebras de tipo finito sobre \(\mathbb Z\)). Sea, pues, \(S=\varprojlim S_\lambda\). Según EGA IV 8.5.2 y 8.5.5, existen un índice \(\lambda\) y un \(\mathcal O_{S_\lambda}\)-módulo localmente libre de tipo finito \(E_\lambda\) tal que \(E=u_\lambda^*(E_\lambda)\) (denotando por \(u_\lambda\) el morfismo \(S\to S_\lambda\), y por \(u_{\lambda\mu}\) los morfismos de transición \(S_\mu\to S_\lambda\)).

Si se pone, para \(\mu\geq\lambda\), \(E_\mu=u_{\lambda\mu}^*(E_\lambda)\), y \(X_\mu=\mathbf P(E_\mu)\), se tiene, para todo \(\mu\), \(X=X_\mu\times_{S_\mu}S\), y \(X=\varprojlim X_\mu\). Por consiguiente, existen un índice \(\lambda'\geq\lambda\) y un haz \(F_{\lambda'}\) de presentación finita sobre \(X_{\lambda'}\), tal que \(F=u'^*_{\lambda'}(F_{\lambda'})\) (siendo \(u'_{\lambda'}\) el morfismo \(X\to X_{\lambda'}\)). Como los \(X_\mu\) y los \(F_\mu\) son de presentación finita, se puede, según EGA IV 11.2.6, encontrar un índice \(\lambda''\) tal que \(F_{\lambda''}\) sea plano sobre \(S_{\lambda''}\). Se verifican, pues, todas las hipótesis de 1.3, lo que concluye la demostración.

